\chapter{Bringing Up the Octal PSRAM}
\label{ch:psram}

Many ESP32-S3 modules package \emph{external} pseudo-static RAM (PSRAM) alongside
the SoC -- but whether there is any, how much, and of what kind depends on the
module. Espressif's S3 family ships in several variants: no PSRAM at all, 2\,MB of
\emph{quad} (4-line) PSRAM, or up to 8\,MB of \emph{octal} (8-line) PSRAM, and from
more than one memory vendor. This chapter is about the octal variant -- the
high-bandwidth option -- which the boards used here happen to populate with an
8\,MB AP~Memory die. The bring-up does not hard-code that, though: it reads the
device's own mode registers and adapts to whichever octal part is fitted (vendor,
density, voltage, latency), so the 8\,MB and the exact constants below are what
\emph{these} parts reported, not universal facts.

Why bother? Internal SRAM is about 292\,KB -- generous for a microcontroller, but
nowhere near enough for the megabyte-scale buffers, heaps, and elaboration stacks
that a full Ada program wants. PSRAM closes that gap, but only if you can turn it
on, and turning it on is one of the most delicate sequences in the whole bring-up:
a high-speed double-data-rate (DDR) octal link that has to be configured register
by register, and then \emph{timed}, before a single byte will round-trip.

This chapter is the deep version of the PSRAM paragraph in the bootloader chapter
(\S\ref{ch:bootloader}). There we noted that the second-stage loader brings PSRAM
up before the application runs; here we tell the whole story -- the deterministic
register script, the empirical timing problem at 80\,MHz, and the de-blobbing
effort that replaced the vendored ESP-IDF binaries with readable from-source code.
It is a genuine war story, and the registers and values below are the real ones,
captured from working silicon.

\section{What the octal PSRAM is, and why it is hard}\index{PSRAM}\index{octal SPI}\index{MSPI}

The PSRAM is an external DRAM die that speaks an \emph{octal} (8-data-line) variant
of the memory-SPI (MSPI) protocol, clocked DDR -- data on both clock edges. The
S3's MSPI controller can map it transparently through the data cache, so once it is
up the chip appears as ordinary memory at \texttt{0x3D000000} and Ada code reads and
writes it with no special API at all (\S\ref{ch:psram-validate}). That transparency
is the whole point, and also the whole danger: a load instruction that targets a
mis-timed PSRAM word does not raise an exception, it \emph{stalls the bus}.

Physically the link uses dedicated MSPI pads: eight data lines (the four ordinary
SPI data pins \texttt{SPID}/\texttt{SPIQ}/\texttt{SPIWP}/\texttt{SPIHD} plus
\texttt{SPID4}--\texttt{SPID7}), a data strobe \texttt{DQS}, the shared
\texttt{SPICLK}, and a second chip select \texttt{CS1} on GPIO\,26. These are the
same pins the hardware chapter (\S\ref{ch:hardware}) flags as off-limits: GPIO\,26
through 37 are wired to the in-package flash and PSRAM, and driving them as general
I/O corrupts the very bus the program is executing from. The PSRAM bring-up is the
code that legitimately owns those pins.

The reason 8 octal DDR data lines at 80\,MHz is hard -- and the reason this chapter
exists -- is timing. At that rate the round-trip from the controller's clock edge
out to the chip and back as data is a meaningful fraction of one bit period, so the
controller cannot just sample data on a fixed edge; it has to \emph{learn} when the
returning bits are valid. That learning step is the only part of the whole sequence
that is not a fixed register script, and it is where almost all the difficulty lives.

\section{The de-blobbing goal}\index{de-blobbing}\index{ESP-IDF}

A working PSRAM bring-up already existed: the vendored ESP-IDF v5.4.4 binaries
\texttt{esp\_psram\_impl\_octal.c.obj} and the \texttt{mspi\_timing\_*.c.obj} set,
linked as opaque blobs. The goal of this effort was to replace those blobs with
readable from-source C (plus a handful of ROM calls), so that nothing in the boot
path is a black box. That was done in two stages, retiring five blobs:

\begin{itemize}
\item \textbf{Stage 1} retired the four timing/GPIO blobs
(\texttt{mspi\_timing\_by\_mspi\_delay}, \texttt{mspi\_timing\_config},
\texttt{mspi\_timing\_tuning}, \texttt{gpio\_periph}), replacing them with
\texttt{mspi\_timing\_src.c}: the clock dividers (core 160\,MHz $\rightarrow$ PSRAM
80\,MHz), the din application, and the pad drive strength.
\item \textbf{Stage 2} retired the last blob,
\texttt{esp\_psram\_impl\_octal.c.obj}, replacing it with
\texttt{psram\_impl\_src.c}: the chip-side mode-register programming and the SPI0
cache-phase configuration.
\end{itemize}

What was \emph{kept} is the ROM: \texttt{esp\_rom\_opiflash\_*} (the octal command
helper and pad config) and \texttt{Cache\_*} (cache disable/enable and the d-bus MMU
map). These are masked into the chip, not vendored blobs, so they cost nothing to
keep and de-ROM-ing them buys little. The net result is the few-dozen-line bring-up
the bootloader chapter advertised.

\cnote{In ESP-IDF this is \texttt{esp\_psram\_impl\_octal} -- thousands of lines of
generated, config-gated code you never see. Here it is two short C files you can
read top to bottom: the deterministic register script, and one honest timing sweep.}

\section{The deterministic part: a fixed register script}\index{register script}\index{mode registers}

Most of the bring-up is a fixed sequence of register writes and ROM calls that is
identical on every boot and every board. The driving function is
\texttt{psram\_bringup()} in \texttt{psram\_boot.c}, which wraps the from-source
enable, maps the chip through the cache, and then tunes the timing:

\begin{lstlisting}[style=sh]
esp_rom_opiflash_pin_config()          # ROM: wire the octal pads SPID4-7 + DQS
mspi_timing_set_pin_drive_strength()   # from-src: FUN_DRV=3 on the 9 octal pads
psram_impl_enable_src()                # from-src: the chip bring-up (below)
psram_impl_get_physical_size_src(&sz)  # from-src: report decoded density -> 8 MB
Cache_Disable_DCache()                 # ROM
Cache_Dbus_MMU_Set(.., 0x3D000000, ..) # ROM: map PSRAM pages x 64 KB
Cache_Enable_DCache(0)                 # ROM
psram_tune_din()                       # from-src: the 80 MHz din sweep (below)
\end{lstlisting}

The pad wiring is \emph{FIX 1}: the ROM flash setup leaves only four of the eight
data lines connected because flash is quad. Without \texttt{SPID4}--\texttt{SPID7}
and \texttt{DQS}, the chip's mode-register read comes back with half its bits
floating, so the density field decodes to \texttt{0x5} (16\,MB) instead of the real
\texttt{0x3} (8\,MB) and the chip is mis-configured from the start.

Inside \texttt{psram\_impl\_enable\_src()} the order is load-bearing: controller
setup, drop to a slow clock, do the chip transactions, raise the clock, then the
cache phases. The deterministic steps are:

\begin{longtable}{p{2.4cm} p{4.2cm} p{5.4cm}}
\toprule
\textbf{Step} & \textbf{Register / call} & \textbf{What it does} \\ \midrule
\endhead
init pins      & \texttt{0x6000906C}=\texttt{0xF00}        & route GPIO\,26 to \texttt{SPICS1}, \texttt{FUN\_DRV}=3 \\
clk pad drive  & \texttt{0x600033FC}=\texttt{0x0210105F}   & SMEM \texttt{SPICLK} pad drive \\
CS timing      & \texttt{0x600030DC}=\texttt{0x0400B18F}   & \texttt{SMEM\_AC}: CS hold 3, setup 3, hold-delay 2 \\
ECC off        & clear \texttt{0x60026058} bit 8           & \texttt{SRAM\_ACE0\_ATTR}: disable inline ECC \\
low speed      & \texttt{enter\_low\_speed\_mode}          & core 80\,MHz, flash+PSRAM /4 = 20\,MHz, clear din \\
var-dummy      & \texttt{0x600020E0} \verb#|=# 2           & SPI1 \texttt{SPI\_FMEM\_VAR\_DUMMY}; no DTR swap \\
write MR0      & \texttt{exec\_cmd} \texttt{0xC0C0}        & \texttt{lt}=1 (fixed), \texttt{read\_latency}=2, \texttt{drive\_str}=0 $\rightarrow$ \texttt{0x28} \\
connect probe  & \texttt{exec\_cmd} \texttt{0x8080}/\texttt{0x0000} & write+read-back \texttt{0x5a6b7c8d} at addr 0 \\
read MR0--8    & \texttt{exec\_cmd} \texttt{0x4040}        & decode vendor / density / Vcc / latency \\
\textbf{tune din} & \texttt{psram\_tune\_din}              & \textbf{empirical 80\,MHz sweep -- the non-deterministic step} \\
high speed     & \texttt{enter\_high\_speed\_mode}         & core 160\,MHz, /2 = 80\,MHz, apply din \\
cache phases   & 5 SPI0 regs (below)                       & program transparent OPI-DDR cache access \\
\bottomrule
\end{longtable}

A few of these deserve a closer look. The MR0 write is a read-modify-write of the
low six bits to \texttt{0x28}: \texttt{lt=1} selects \emph{fixed} read latency (so
the controller knows exactly how many dummy cycles to insert), and
\texttt{read\_latency=2} encodes a 10-cycle fixed latency. The connectivity probe is
the only ``is it there?'' test in the octal path -- there is no JEDEC \texttt{0x9F}
read-ID -- so it writes the witness word \texttt{0x5a6b7c8d} to address 0 and reads
it back; identity comes entirely from the mode registers. On a genuine
ESP32-S3-DevKitC the decode prints:

\begin{lstlisting}[style=sh]
PSRAM: AP Memory octal DDR @80MHz, 64 Mbit (8 MB), dev gen 3, Vcc 3.0V
  latency: read 10-cyc (fixed), write 5-cyc;  MR0=28 MR1=0d MR2=93 MR3=60 MR4=40 MR8=05
\end{lstlisting}

All the mode-register traffic -- the MR write, the probe, the MR reads -- goes
through the ROM helper \texttt{esp\_rom\_opiflash\_exec\_cmd} in OPI-DTR mode
(\texttt{mode 7}) with a 16-bit command, 32-bit address, and the PSRAM chip select
\texttt{CS=BIT(1)}. Read command is \texttt{0x4040}, write \texttt{0xC0C0}, sync read
\texttt{0x0000}, sync write \texttt{0x8080}; read dummy 18 half-cycles, write dummy 8.

\subsection{The golden cache-phase registers}\index{cache MMU}

The last deterministic step configures the SPI0 (cache) controller so that a plain
load from \texttt{0x3D000000} silently becomes an octal-DDR PSRAM read. These five
register values were captured from the working blob over JTAG -- the ``golden''
state -- and are written directly. They encode the same OPI-DDR geometry as the
manual transactions above, but for the transparent cache path:

\begin{longtable}{p{2.6cm} p{2.4cm} p{6.8cm}}
\toprule
\textbf{Register} & \textbf{Value} & \textbf{Meaning} \\ \midrule
\endhead
\texttt{SRAM\_CMD}  & \texttt{0x007C0000} & octal command/address/data lanes \\
\texttt{SRAM\_DRD\_CMD} & \texttt{0xF0000000} & cache read command \texttt{0x0000}, 16-bit \\
\texttt{SRAM\_DWR\_CMD} & \texttt{0xF0008080} & cache write command \texttt{0x8080}, 16-bit \\
\texttt{SMEM\_DDR}  & \texttt{0x00003023} & DDR enable + variable dummy \\
\texttt{CACHE\_SCTRL} & \texttt{0x01F7C479} & 32-bit addr, read dummy 18 / write 8, octal lanes \\
\bottomrule
\end{longtable}

\texttt{CACHE\_SCTRL} carries the lane bits that make every phase octal
(\texttt{SCMD\_OCT}, \texttt{SADDR\_OCT}, \texttt{SDOUT\_OCT}, \texttt{SDIN\_OCT},
\texttt{SRAM\_OCT}) and the \texttt{SCMD\_4BYTE} flag for the 32-bit address. After
these five writes, \texttt{Cache\_Resume\_DCache(0)} hands the cache back and the
chip is, deterministically, configured. Everything so far is reproducible from
source. What remains is the one step that is not.

\section{The non-deterministic part: timing tuning}\index{din timing}

At 80\,MHz DDR the controller must choose \emph{when} to sample the returning data.
The S3 exposes three per-data-line knobs for this, all in the SPI0 (cache) MSPI block
at \texttt{0x6000\_30xx}:

\begin{itemize}
\item \texttt{din\_mode} -- \texttt{SMEM\_DIN[0..7]\_MODE} (the register at
\texttt{0x600030C0}): the sampling edge/phase, 3 bits per lane, 8 settings.
\item \texttt{din\_num} -- \texttt{SMEM\_DIN[0..7]\_NUM} at \texttt{0x600030C4}: an
integer-cycle delay added to data-in.
\item \texttt{extra\_dummy} -- \texttt{SMEM\_TIMING\_CALI} at \texttt{0x600030BC}:
extra dummy cycles before the data window opens.
\end{itemize}

Of these, \texttt{din\_mode} is the one that matters here, and the one whose correct
value is not knowable from source -- it depends on the actual chip, controller, and
temperature. Picking it wrong does not fail gracefully: a load through the cache at a
bad \texttt{din\_mode} stalls the bus with no exception, the PC pins, and no interrupt
can preempt it. So this knob has to be chosen carefully, and chosen safely.

\section{The tuning war story}

\subsection{Why IDF's own tuning is a no-op}\index{din\_mode}

The natural assumption is that the vendor blob already calibrates \texttt{din\_mode}.
It does not. Single-stepping \texttt{mspi\_timing\_psram\_tuning} on live silicon
showed why: the IDF sweep drops the MSPI to \emph{20\,MHz} before testing and only
restores 80\,MHz afterwards. At 20\,MHz the data eye is enormous, so the sampling
phase is irrelevant -- every one of the 14 candidate configurations reads the
reference pattern back byte-perfectly. The DTR selection heuristic sees an
implausibly wide all-pass window (run length $\ge$ 6) and \emph{distrusts} it, falling
back to \texttt{default\_config\_id = 5}, which is \texttt{din\_mode = 4}.

\cnote{The captured ``calibration'' is really a vendor-default lookup. In the blob,
\texttt{din\_mode = 4} is dead-centre of the \emph{failing} region at 80\,MHz -- it
crashes the cache read. That is why the old shim carried a hand-coded ``FIX 2''
forcing \texttt{din\_mode = 1}: the only way to make the blob's chosen default usable
was to override it.}

So the working configuration was, historically, ``vendor default plus a magic
constant.'' De-blobbing meant either reproducing that magic constant honestly, or
measuring the right value for real.

\subsection{Why you cannot naively sweep the cache path}

The obvious fix is to sweep \texttt{din\_mode} yourself: set each mode, read a known
pattern back through the cache, keep the modes that pass. This was tried, and the
hardware fights back at every turn:

\begin{enumerate}
\item \textbf{A wrong cache din hangs the CPU.} A synchronous PSRAM load with a bad
din never completes -- the bus stalls, the PC pins, \texttt{exccause=0}, and no
interrupt can preempt a stalled load. You cannot ``read it and check if it is
garbage,'' because the read does not return.
\item \textbf{Async preload dodges the hang but poisons the controller.} Driving the
sweep with \texttt{Cache\_Start\_DCache\_Preload} plus a polled timeout lets the CPU
survive (it only polls a flag), but a stalled preload leaves the MSPI/cache
controller wedged: subsequent cache reads return the fault-fill pattern
\texttt{0xbad00bad}, and \emph{even \texttt{Cache\_Disable\_DCache} hangs}. Nothing
short of an MSPI peripheral reset clears it -- which also kills flash execute-in-place.
\item \textbf{SPI1 has its own din registers.} SPI1 (the manual-transaction path) keeps
separate din registers at \texttt{0x6000\_20xx}, distinct from the cache controller's
\texttt{0x6000\_30xx}. A bounded SPI1 read calibrates the wrong register -- it tells
you nothing directly about the cache path.
\end{enumerate}

The conclusion is uncomfortable but precise: the din-sensitive path \emph{is} the
cache path, and the cache path cannot be safely probed because any bad value risks an
unrecoverable wedge.

\subsection{The fix: a bounded SPI1 sweep at the real speed}

\texttt{psram\_tune\_din()} threads the needle. It sweeps \texttt{din\_mode} at the
real \emph{80\,MHz} operating clock, but it does the readback over a \emph{bounded SPI1
manual transaction} (\texttt{esp\_rom\_opiflash\_exec\_cmd}), not through the cache. A
wrong din over SPI1 returns garbage and the transaction \emph{still completes} -- so
the sweep can measure all eight modes without ever stalling the bus. SPI0 and SPI1 share
the same MSPI clock and pads, so the window measured over SPI1 is, to a known offset,
the window the cache wants. The core of it:

\begin{lstlisting}[style=sh]
# write a reference word (data-out: din-independent), then sweep the read din:
exec_cmd(spi1, OPI_DTR, 0x8080, addr, dummy=8,  &ref)      # reference write
for m in 0..7:
    REG(0x600030C0) = din_word(m)   # SPI0 (cache) din-mode, 3 bits x 9 lanes
    REG(0x600020C0) = din_word(m)   # SPI1 (manual) din-mode
    exec_cmd(spi1, OPI_DTR, 0x0000, addr, dummy=18, &rd)   # bounded read
    if rd == ref: pass |= 1 << m
\end{lstlisting}

The measured window is \texttt{passmask = 0xC3} -- modes \{6, 7, 0, 1\} pass, modes
\{2, 3, 4, 5\} fail. Crucially this is a \emph{circular} arc: \texttt{din\_mode} is a
sampling phase, so mode 7 is adjacent to mode 0, and the passing run wraps the
$7\rightarrow0$ seam. The selection logic walks the modes twice (two laps) to catch a
window that wraps, and picks the centre of the widest passing run:

\begin{lstlisting}[style=sh]
mode 0,1 PASS | mode 2,3 fail (0x..76) | mode 4,5 fail (0x..e2) | mode 6,7 PASS
                                                       passmask 0xC3  ->  centre = mode 0
\end{lstlisting}

Notice that IDF's default \texttt{din\_mode = 4} sits dead-centre in the \emph{failing}
region -- the 80\,MHz sweep correctly rejects exactly what the 20\,MHz tune handed
back. The measured centre (mode 0) has more margin than FIX 2's mode 1, which was a
valid but edge-of-eye point. If the sweep is degenerate (empty or all-pass) the code
falls back to mode 1. (One subtlety, recorded in the source: the SPI1-measured centre
can sit slightly off the cache path's true optimum, and that offset is invisible to a
quick checksum but can corrupt a heap-heavy app under sustained access; so when mode 1
-- the validated-robust cache point for this chip -- lies inside the measured window,
the code prefers it.)

Because din timing affects \emph{reads} only and the bootloader itself runs from
flash/IRAM rather than PSRAM, retuning the PSRAM din live, after the chip is mapped, is
safe -- nothing the bootloader is executing depends on it.

\subsection{Cross-validation on two boards}

The strongest evidence that this measures a real property and not a per-board fluke is
that two physically different boards land on the \emph{same} window. The whole effort
was developed on a Meshnology LoRa AIOT board; the finished from-source bring-up was
then run unchanged on a genuine Espressif ESP32-S3-DevKitC. Both independently measured
\texttt{passmask = 0xC3} $\rightarrow$ mode 0, on the same AP Memory 8\,MB octal die.
The din window tracks a stable chip-plus-controller property, which is exactly what you
want from a calibration: reproducible, not lucky.

\section{The 80\,MHz versus 40\,MHz trade}

It is worth being explicit about why this complexity exists at all because there is a
simpler road not taken. At 40\,MHz the MSPI does \emph{not} need tuning: below the
40\,MHz threshold the data eye is wide enough that the sampling phase is irrelevant,
din/num/dummy are all zero, and the entire bring-up collapses into a straight register
script with no empirical step -- every value a compile-time constant. That is the
clean, fully-deterministic option.

The cost is bandwidth: 40\,MHz DDR is half the throughput of 80\,MHz DDR. This build
chose 80\,MHz, and therefore inherited the tuning problem because the PSRAM here is
used for bandwidth-sensitive work, not just capacity. The 40\,MHz script remains a
documented fallback -- the right choice if a future board needs PSRAM only as a large,
slow backing store and would rather not carry a sweep at all.

\section{On-target validation}\index{external RAM}
\label{ch:psram-validate}

None of this matters unless an Ada program can actually use the result, so the
\texttt{esp32s3\_psram} example exists to prove it does. It places a 1\,MB static byte
array in external RAM -- no heap, no \texttt{malloc} -- using a linker section that
\texttt{psram.ld} maps onto the PSRAM region:

\begin{lstlisting}
Buffer : array (0 .. 1024 * 1024 - 1) of Unsigned_8
  with Linker_Section => ".ext_ram.bss";
\end{lstlisting}

\texttt{Big.Run} fills the buffer with a known pattern (byte $i$ gets $i \bmod 256$),
reads every byte back, and checksums the lot. Because the pattern is bytes
$0\ldots255$ repeating across exactly $2^{20}$ bytes, the sum is a fixed,
hand-computable constant: \texttt{checksum = 0x07F80000}. Any timing error anywhere in
the 1\,MB round-trip through the cache would perturb it. On hardware:

\begin{lstlisting}[style=sh]
[ada-free-boot] octal PSRAM up: rc=0  8 MB
[ada-free-boot] PSRAM mapped @0x3D000000 rc=0
[ada-free-boot] PSRAM din tuned @80MHz: passmask=0xc3 -> mode 0
[psram] buffer @ 0x3d000000  1048576 bytes  checksum=0x07f80000  (PSRAM)
\end{lstlisting}

The buffer's \texttt{0x3D000000} address confirms it really lives in the external-RAM
range, and the deterministic checksum -- byte-identical across resets and across both
test boards -- confirms the full megabyte round-trips correctly through a cache served
by a din the bootloader measured for itself. There is no FreeRTOS, no ESP-IDF, and no
remaining blob in the path: the entire octal-PSRAM bring-up is now readable source plus
a few ROM calls, which is exactly what de-blobbing set out to achieve.

One caveat the example notes for completeness: task stacks stay in internal SRAM. PSRAM
is unreachable whenever the cache is disabled -- for instance during a flash write --
so it is the right home for large data buffers and heaps, but not for a stack that an
interrupt might unwind while the cache is down (\S\ref{ch:bootloader}).
